% theme: numerical_calculation
\MajorQuestionLesson{数値計算}
\SetRowSpace{3.5mm}

\TypeQuestionSingle{次の小数を、既約分数で表しなさい。}{shousuu_bunsuu}
\QQRow{$12.5$}{}{$0.025$}{}
\QQRow{$0.125$}{}{$0.75$}{}
\QQRow{$1.25$}{}{$0.32$}{}
\QQRow{$0.875$}{}{$0.45$}{}

\TypeQuestionSingle{筆算を使わず、小数を分数に直し次の計算をしなさい。}{keisan_kufuu}
\QQRow{$12.5\times1.2$}{}{$2.5\times3.6$}{}
\QQRow{$0.75\times1.6$}{}{$1.25\times4.8$}{}
\QQRow{$0.24\times25$}{}{$0.125\times56$}{}
\QQRow{$6.25\times3.2$}{}{$0.48\times125$}{}
%割り算をあと２つ増やして下さい。答えは分数にならない(小数で表せる)ように。
\QQRow{$3.2\div0.8$}{}{$7.5\div1.25$}{}
\QQRow{$0.84\div0.35$}{}{$4.8\div0.16$}{}

\TypeQuestionSingle{約分・計算順序を工夫して、次の計算をしなさい。}{junjo}
%割り算も混ぜて下さい。答えは分数にならない(小数で表せる)ように。
\QQRow{$2.5\times4\times3.7$}{}{$1.25\times8\times4.6$}{}
\QQRow{$0.4\times7.5\times5$}{}{$0.25\times1.6\times12$}{}
\QQRow{$6.4\times0.125\times2.5$}{}{$3.6\times2.5\times0.4$}{}
\QQRow{$12.5\times0.8\times0.25$}{}{$0.32\times25\times1.25$}{}

\LessonPageBreak

\TypeQuestionSingle{約分・計算順序を工夫して、次の計算をしなさい。}{yakubun}
%答えは分数にならない(小数で表せる)ように。問題数は変えないで
\QQRow{$\dfrac{2.4\times15}{0.6\times8}$}{}{$\dfrac{3.5\times1.2}{0.7\times0.8}$}{}
\QQRow{$\dfrac{0.45\times14}{2.1\times0.3}$}{}{$\dfrac{1.25\times3.6}{0.75\times2.5}$}{}
\QQRow{$\dfrac{4.8\times0.35}{0.56\times1.5}$}{}{$\dfrac{0.72\times2.5}{0.15\times4}$}{}
\QQRow{$\dfrac{2.5\times3.6\times0.8}{1.2\times0.5}$}{}{$\dfrac{0.48\times7.5}{1.25\times0.24}$}{}

\TypeQuestionSingle{分配法則を利用して、次の計算をしなさい。}{bunpai}
%下記タイプの分配法則活用が圧倒的に有効なので、これをあと４題追加して。
\QQRow{$12.5\times1.6+7.5\times1.6$}{}{$4.8\times2.5-0.8\times2.5$}{}
\QQRow{$\dfrac{3.6\times5+1.4\times5}{2.5}$}{}{$\dfrac{8.4\times1.25-0.4\times1.25}{0.5}$}{}

\TypeQuestionSingle{各式に与えられた値を代入し、未知変数の値を求めなさい。}{koushiki_dainyuu}
%6題追加して下さい。このタイプは実践的で大事！
%文字に添え字があるタイプも追加して
%右辺に代入して左辺をもとめるのではなく両辺に代入して方程式を解くようなものも追加。
%v=v_0+atに、v,a,v_0を入れる、みたいな。
\QQ{$v=\dfrac{x}{t}$ に、$x=12.5$、$t=0.25$ を代入する。}{}
\QQ{$F=ma$ に、$m=1.25$、$a=4.8$ を代入する。}{}
\QQ{$E=\dfrac{1}{2}mv^2$ に、$m=2.5$、$v=4.8$ を代入する。}{}
\QQ{$Q=mc\Delta T$ に、$m=0.25$、$c=4.2$、$\Delta T=12$ を代入する。}{}
